\title{Group Sequence Policy Optimization}

\begin{abstract}
We present Group Sequence Policy Optimization (GSPO), a robust and effective reinforcement learning method developed for training large language models. In contrast to prior approaches that rely on importance ratios computed at the token level, GSPO employs a sequence-level perspective—constructing importance ratios, applying clipping, and implementing rewards and optimization procedures at the level of entire sequences. Our experiments show that GSPO not only outperforms the GRPO algorithm in terms of training effectiveness and results but also offers enhanced stability during Mixture-of-Experts (MoE) RL training and promises to streamline RL infrastructure development. The advantages offered by GSPO have played a key role in enabling the significant performance gains of the recent Qwen3 models.
\end{abstract}

\section{Introduction}
\label{sec:intro}

Reinforcement learning (RL) has become a key approach in advancing the scalability of language models \citep{o1,dpsk-r1,qwq32b,qwen3}. Utilizing RL at scale enables language models to acquire the ability to address complex challenges—including high-level mathematics and programming—by engaging in extended and more intricate reasoning tasks.

Achieving effective scaling of RL through increased computational resources critically depends on ensuring training stability and robustness. Nonetheless, leading RL methods such as GRPO \citep{grpo} are plagued by significant instability during the training of extremely large language models, which frequently leads to catastrophic and irrecoverable collapse of the models \citep{qwen3, minimax-m1}. Such instability presents a significant obstacle to advancing the performance of language models by means of extended RL training.

In this study, we demonstrate that GRPO's instability originates from a core flaw in its use of importance sampling weights, which are both incorrectly applied and unreliable within its framework. This misstep results in significant variance in training, leading to accumulating noise as response sequences become longer. Additionally, the clipping mechanism exacerbates this effect, accelerating the process toward model collapse.

In response to these fundamental challenges, we introduce \textbf{Group Sequence Policy Optimization (GSPO)}, an RL framework tailored for optimizing large language models. The primary advancement of GSPO is its rigorously justified importance ratio, formulated in terms of sequence likelihood as referenced by \citep{click}, thus adhering to the foundational concept of importance sampling. Moreover, GSPO derives normalized rewards by evaluating the advantages across multiple generated responses for each query, thereby harmonizing reward allocation at the sequence level with the objectives of optimization.

The results of our empirical assessment reveal that GSPO markedly surpasses GRPO with regard to stability, training efficiency, and overall effectiveness. Importantly, GSPO intrinsically addresses the instability issues often encountered during RL training of large-scale Mixture-of-Experts (MoE) models, removing reliance on elaborate stabilization mechanisms and suggesting avenues for streamlining RL architectures. These advantages have played a pivotal role in realizing the notable performance gains seen in the most recent Qwen3 models. Looking ahead, we consider GSPO to offer a dependable and extensible framework capable of fostering ongoing progress in large-scale RL training for language models.

\section{Preliminaries}

\paragraph{Notation}
Throughout this work, we interpret an autoregressive language model, parameterized by $\theta$, as a policy $\pi_\theta$. The symbol $x$ represents a query, while $\mathcal{D}$ signifies the set of all queries. When considering a response $y$ given query $x$, its probability according to the policy $\pi_\theta$ is expressed as $\pi_\theta(y|x) = \prod_{t=1}^{|y|} \pi_\theta(y_t \mid x, y_{<t})$, where $|y|$ denotes the total number of tokens in $y$. To assess a query-response pair $(x, y)$, we employ a verifier $r$, which assigns a reward $r(x, y)$ in the interval $[0, 1]$.

\paragraph{Proximal Policy Optimization (PPO)} Proximal Policy Optimization (PPO) \citep{ppo} leverages experience collected under the previous policy $\pi_{\theta_\text{old}}$ to perform policy updates that are restricted to a neighborhood of the existing policy, utilizing a clipping technique to enforce this constraint. In particular, PPO optimizes the following objective (excluding the KL regularization term hereafter for clarity, since it is not central to this discussion):

\begin{align}
\mathcal{J}_\text{PPO}(\theta) = \mathbb{E}_{ x \sim \mathcal{D},\, y \sim \pi_{\theta_\text{old}}( \cdot | x) }
\left[ \frac{1}{|y|} \sum_{t=1}^{|y|} 
\min \left( w_{t}(\theta) \widehat{A}_{t},  \, \mathrm{clip} \left( w_{t}(\theta), 1 - {\varepsilon}, 1 + {\varepsilon}\right) \widehat{A}_{t} \right)
\right],
\end{align}

Here, the importance ratio assigned to the token $y_t$ is given by $
w_{t}(\theta) = \frac{ \pi_{\theta} (y_{t} | x, y_{<t}) }{ \pi_{\theta_\text{old}} (y_{t} | x,y_{<t})} $. The value of the advantage $\widehat{A}_{t}$ for $y_t$ is calculated using a separate value estimator, and $\varepsilon$ specifies the interval within which the importance ratios are clipped.

A fundamental obstacle with PPO in real-world settings is its strong dependence on the value model. In particular, since the value model is typically comparable in size to the policy model, this results in significant demands on both memory and computational resources. Moreover, the success of the algorithm is closely tied to the trustworthiness of the value estimates it provides. Obtaining a robust value model is already nontrivial; extending its effectiveness to accommodate longer outputs and increasingly complex tasks exacerbates these difficulties even further.

\paragraph{Group Relative Policy Optimization (GRPO)} GRPO \citep{grpo} eliminates the dependency on a value network by determining the relative advantage among a set of candidate responses provided for an identical prompt. More precisely, GRPO maximizes the following objective:

\begin{align}
\label{equ:grpo}
\mathcal{J}_\text{GRPO}(\theta) = \mathbb{E}_{ x \sim \mathcal{D},\, \{y_i\}_{i=1}^G \sim \pi_{\theta_\text{old}}( \cdot | x) }
\left[ \frac{1}{G} \sum_{i=1}^{G} \frac{1}{|y_i|} \sum_{t=1}^{|y_i|} 
\min \bigg( w_{i,t}(\theta) \, \widehat{A}_{i,t},\;  \mathrm{clip} \big( w_{i,t}(\theta), 1 - {\varepsilon}, 1 + {\varepsilon}\big) \, \widehat{A}_{i,t} \bigg)
\right],
\end{align}

where $G$ is the number of generated responses to each query $x$ (i.e., the group size), and the importance ratio $w_{i,t}(\theta)$ and advantage $\widehat{A}_{i,t}$ of token $y_{i,t}$ are:

\begin{align}
    w_{i,t}(\theta)=\frac{ \pi_{\theta} (y_{i,t} | x, y_{i,<t}) }{ \pi_{\theta_\text{old}} (y_{i,t} | x,y_{i,<t})},\quad\ 
    \widehat{A}_{i,t} = \widehat{A}_{i} = \frac{r(x, y_i) - \mathrm{mean} \left( \{ r(x, y_i) \}_{i=1}^G \right) }{ \mathrm{std} \left( \{ r(x, y_i) \}_{i=1}^G \right) },
\end{align}

Here, $w_{i,t}(\theta)$ denotes the importance weight for the $t$-th token of sample $i$ under the new model parameters $\theta$ relative to the prior parameters $\theta_\text{old}$. The advantage estimate $\widehat{A}_{i,t}$, which is identical across all time steps for a given sample $i$, is defined as the standardized reward for that sample, computed by subtracting the mean reward over the batch $\{ r(x, y_i) \}_{i=1}^G$ and normalizing by its standard deviation.

respectively, where all the tokens in $y_i$ share the same advantage as $\widehat{A}_{i}$.

\section{Motivation}
\label{sec:motivation}

As models become larger, incorporate increased sparsity (as seen in Mixture-of-Experts architectures), and generate longer responses, the need for substantial rollout batch sizes becomes more pronounced in order to fully leverage available hardware during reinforcement learning. In practice, to make more efficient use of collected samples, these large rollout batches are typically subdivided into several mini-batches, each used for individual gradient updates. However, this strategy naturally results in an off-policy learning scenario: the response samples $y$ originate from a previously used policy, $\pi_{\theta_\text{old}}$, instead of the currently updated policy $\pi_\theta$. Consequently, mechanisms such as clipping in PPO and GRPO become essential to restrict the influence of samples that deviate too far from the current policy during gradient computation.

Although techniques such as clipping are employed to address off-policy discrepancies, we observe that GRPO suffers from a deeper problem: \textit{the objective itself lacks well-posedness}. This challenge is especially pronounced when training large-scale models on tasks requiring extended responses, often resulting in severe model degeneration. The root of this issue in the GRPO objective lies in the incorrect application of importance sampling weights. Importance sampling is designed to compute expectations of a function $f$ with respect to a target distribution $\pi_\text{tar}$ by appropriately adjusting the weights of samples generated from a behavior distribution $\pi_\text{beh}$:

\begin{align}
\mathbb{E}_{z \sim \pi_\text{tar}} \left[ f(z) \right] = \mathbb{E}_{z \sim \pi_\text{beh}} \left[ \frac{ \pi_\text{tar} (z) }{ \pi_\text{beh} (z)} \, f(z) \right].
\end{align}

A key aspect of this approach is that it requires taking an average over many samples ($N \gg 1$) drawn from the behavior distribution $\pi_\text{beh}$ so that the importance weight $\frac{ \pi_\text{tar} (z) }{ \pi_\text{beh} (z)}$ can adequately address the discrepancy between distributions.

On the other hand, GRPO introduces the importance weight $\frac{ \pi_{\theta} (y_{i,t} | x, y_{i,<t}) }{ \pi_{\theta_\text{old}} (y_{i,t} | x,y_{i,<t})}$ at every individual token position $t$. This weighting relies solely on a particular sampled token $y_{i,t}$ obtained from the distribution $\pi_{\theta_\text{old}}( \cdot | x, y_{i, <t})$ for each timestep, and as a result, does not accomplish the expected correction of the distribution. Rather, this approach significantly amplifies the variance of the gradients during training, particularly for lengthy sequences, and this variance is further aggravated by the use of clipping. Through empirical studies, we found that such instability may provoke a collapse of the model, from which recovery is often unattainable. Once such a failure manifests, restarting training—even when rolling back to earlier checkpoints, extensively optimizing hyperparameters (such as clipping thresholds), increasing the context length, or altering the reinforcement learning queries—proves ineffective.

This observation highlights an underlying flaw in the architecture of GRPO. Specifically, the ineffectiveness of token-level importance weighting underscores an essential concept: \textbf{the granularity of the optimization target must correspond to that of the reward}. As the reward is assigned at the sequence level, implementing off-policy corrections at the token level introduces inconsistencies. Therefore, we are driven to abandon the token-level formulation in favor of leveraging importance weighting and conducting optimization processes at the \textit{sequence level} instead.

\section{Algorithm}

\subsection{GSPO: Group Sequence Policy Optimization}

Although the token-level importance weight $\frac{ \pi_{\theta} (y_{i,t} | x, y_{i,<t}) }{ \pi_{\theta_\text{old}} (y_{i,t} | x,y_{i,<t})}$ poses challenges in the GRPO framework, our analysis reveals that, specifically for language generation, the \textit{sequence-level} importance weight $\frac{ \pi_{\theta} (y | x) }{ \pi_{\theta_\text{old}} (y | x)}$ possesses a well-defined theoretical interpretation: it quantifies the extent to which a response $y$ drawn from $\pi_{\theta_\text{old}} (\cdot | x)$ diverges from the target distribution $\pi_{\theta} (\cdot | x)$. This relationship naturally corresponds to rewards measured at the sequence level and provides a relevant signal for guiding the clipping mechanism.

Based on this straightforward observation, we propose the \textbf{Group Sequence Policy Optimization (GSPO)} algorithm. GSPO employs the following sequence-level optimization objective:

\begin{align}
\mathcal{J}_\text{GSPO} (\theta) =
\mathbb{E}_{ x \sim \mathcal{D},\, \{y_i\}_{i=1}^G \sim \pi_{\theta_\text{old}}( \cdot | x) }
\left[ 
\frac{1}{G} \sum_{i=1}^{G}
\min \left( s_{i}(\theta)  \widehat{A}_{i},  \, \mathrm{clip} \left( s_{i}(\theta), 1 - {\varepsilon}, 1 + {\varepsilon} \right) \widehat{A}_{i} \right) 
\right],
\label{equ:gspo}
\end{align}

where we adopt the group-based advantage estimation:

\begin{align}
\widehat{A}_{i} = \frac{r(x, y_i) - \mathrm{mean} \left( \{ r(x, y_i) \}_{i=1}^G \right) }{ \mathrm{std} \left( \{ r(x, y_i) \}_{i=1}^G \right) },
\end{align}

and define the importance ratio $s_{i}(\theta)$ based on sequence likelihood \citep{click}:

\begin{align}
s_{i}(\theta) = \left( \frac{ \pi_{\theta} (y_i | x) }{ \pi_{\theta_\text{old}} (y_i | x)} \right)^{\frac{1}{|y_i|}}
=
\exp \left( \frac{1}{|y_i|} \sum_{t=1}^{|y_i|} \log \frac{ \pi_{\theta} (y_{i,t} | x, y_{i,<t}) }{ \pi_{\theta_\text{old}} (y_{i,t} | x,y_{i,<t})} \right).
\end{align}

Alternatively, $s_i(\theta)$ can be rewritten in two equivalent forms: The first expression represents the $|y_i|$-th root of the ratio between the likelihoods under the current and previous parameterizations, $\theta$ and $\theta_\text{old}$, respectively. The second expression reformulates this by exponentiating the mean of the log-likelihood ratios across each token $t$ in the sequence, reflecting the tokenwise average log-ratio between the policies over the trajectory $y_i$ conditioned on $x$.

Consequently, GSPO implements clipping at the level of entire responses rather than at individual tokens. This approach effectively filters out highly ``off-policy'' samples during gradient computation, thereby aligning with both the reward assessment and optimization processes that occur at the sequence level. We additionally employ length normalization in $s_{i}(\theta)$, which serves to minimize variance and maintain $s_{i}(\theta)$ within a consistent numerical interval. In the absence of such normalization, significant modifications in the likelihood of a small number of tokens could lead to excessive variability in the overall sequence-level importance ratio, making it necessary to adopt differing clipping thresholds for responses of varying lengths. It is further important to recognize that the clipping thresholds utilized by GSPO generally differ in scale from those used in earlier methods (such as GRPO), owing to the distinct definitions given to the importance ratios.

\subsection{Gradient Analysis}
\label{subsec:gradient}

We can derive the gradient of the GSPO objective as follows (clipping is omitted for brevity):

\begin{align}
\nabla_{\theta} \mathcal{J}_\text{GSPO} (\theta)
=&\ 
\nabla_{\theta} \mathbb{E}_{ x \sim \mathcal{D},\, \{y_i\}_{i=1}^G \sim \pi_{\theta_\text{old}}( \cdot | x) }
\left[ 
\frac{1}{G} \sum_{i=1}^{G} s_{i}(\theta) \widehat{A}_{i}
\right] \\
= &\ 
\mathbb{E}_{ x \sim \mathcal{D},\, \{y_i\}_{i=1}^G \sim \pi_{\theta_\text{old}}( \cdot | x) }
\left[ 
\frac{1}{G} \sum_{i=1}^{G}
s_{i}(\theta) \widehat{A}_{i}
\cdot \nabla_{\theta} \log s_{i}(\theta)
\right] \\
=&\ 
\mathbb{E}_{ x \sim \mathcal{D},\, \{y_i\}_{i=1}^G \sim \pi_{\theta_\text{old}}( \cdot | x) }
\left[ 
\frac{1}{G} \sum_{i=1}^{G} 
\left( \frac{ \pi_{\theta} (y_i | x) }{ \pi_{\theta_\text{old}} (y_i | x)} \right)^{\frac{1}{|y_i|}} \widehat{A}_{i} 
\cdot \frac{1}{|y_i|} \sum_{t=1}^{|y_i|} \nabla_{\theta} \log \pi_{\theta} (y_{i,t} | x, y_{i,<t}) 
\right].
\label{equ:gspo_grad}
\end{align}

For comparison, the gradient of the GRPO objective is as follows (note that $\widehat{A}_{i,t} = \widehat{A}_{i}$):

\begin{align}
\nabla_{\theta} \mathcal{J}_\text{GRPO} (\theta) 
=&\ 
\nabla_{\theta} \mathbb{E}_{ x \sim \mathcal{D},\, \{y_i\}_{i=1}^G \sim \pi_{\theta_\text{old}}( \cdot | x) }
\left[ \frac{1}{G} \sum_{i=1}^{G} \frac{1}{|y_i|} \sum_{t=1}^{|y_i|} 
w_{i,t}(\theta) \widehat{A}_{i,t}
\right] \\
=&\
\mathbb{E}_{ x \sim \mathcal{D},\, \{y_i\}_{i=1}^G \sim \pi_{\theta_\text{old}}( \cdot | x) }
\left[ \frac{1}{G} \sum_{i=1}^{G} \widehat{A}_{i}
\cdot \frac{1}{|y_i|} \sum_{t=1}^{|y_i|}
\frac{ \pi_{\theta} (y_{i,t} | x, y_{i,<t}) }{ \pi_{\theta_\text{old}} (y_{i,t} | x,y_{i,<t})}
\nabla_{\theta} \log \pi_{\theta} (y_{i,t} | x, y_{i,<t})  
\right].
\end{align}

Accordingly, the principal difference between GSPO and GRPO is in \textit{the manner in which they assign weights to the gradients of the tokens' log likelihoods}. GRPO determines the weight of each token based on its individual ``importance weight'' $\frac{ \pi_{\theta} (y_{i,t} | x, y_{i,<t}) }{ \pi_{\theta_\text{old}} (y_{i,t} | x,y_{i,<t})}$. These variable weights, which can range over $(0, 1+\varepsilon]$ when $\widehat{A}_{i} >0$, or $[1-\varepsilon, +\infty)$ when $\widehat{A}_{i} < 0$, are nontrivial and, through accumulation, may produce erratic effects during the course of training. Conversely, GSPO applies uniform weighting to all tokens within a response, thereby removing this particular source of instability present in GRPO.

\subsection{GSPO-token: A Token-level Objective Variant}

In situations such as multi-turn RL, a more detailed adjustment of advantages beyond the sequence level may be necessary. Therefore, we propose a variant of the GSPO objective referred to as \textbf{GSPO-token}, which enables the modification of advantages on a token-by-token basis:

\begin{align}
\mathcal{J}_\text{GSPO-token}(\theta)
=
\mathbb{E}_{ x \sim \mathcal{D},\, \{y_i\}_{i=1}^G \sim \pi_{\theta_\text{old}}( \cdot | x) }
\left[ 
\frac{1}{G} \sum_{i=1}^{G} \frac{1}{|y_i|} \sum_{t=1}^{|y_i|} 
\min \left( s_{i,t}(\theta) \widehat{A}_{i,t},  \, \mathrm{clip} \left( s_{i,t}(\theta), 1 - {\varepsilon}, 1 + {\varepsilon} \right) \widehat{A}_{i,t} \right) 
\right],
\label{equ:gspo-token}
\end{align}

The expression above defines the objective function for the GSPO-token approach, where for each context $x$ sampled from the dataset $\mathcal{D}$, and each group of $G$ sequences $\{y_i\}$ generated independently by the previous policy $\pi_{\theta_\text{old}}(\cdot|x)$, the expectation is taken. For each sequence $y_i$, the mean is calculated over its tokens, and then across all $G$ sampled sequences. At the token level, the minimum is computed between the product $s_{i,t}(\theta)\widehat{A}_{i,t}$ and the clipped term $\mathrm{clip}(s_{i,t}(\theta), 1 - \varepsilon, 1 + \varepsilon)\widehat{A}_{i,t}$. This structure ensures per-token policy updates are both robust and stabilized by clipping, averaging across tokens and sequences for the overall objective.

where

\begin{align}
s_{i,t}(\theta) 
= \mathrm{sg} \left[ s_{i}(\theta) \right]  \cdot \frac{ \pi_{\theta} (y_{i,t} | x, y_{i,<t}) }{ \mathrm{sg} \left[ \pi_{\theta} (y_{i,t} | x, y_{i,<t}) \right] },
\end{align}

and $\mathrm{sg}[\cdot]$ denotes only taking the numerical value but stopping the gradient, corresponding to the \texttt{detach} operation in PyTorch. The gradient of GSPO-token can be derived as:

\begin{align}
\nabla_{\theta} \mathcal{J}_\text{GSPO-token}(\theta)
=&\ 
\nabla_{\theta} \mathbb{E}_{ x \sim \mathcal{D},\, \{y_i\}_{i=1}^G \sim \pi_{\theta_\text{old}}( \cdot | x) }
\left[ 
\frac{1}{G} \sum_{i=1}^{G}
\frac{1}{|y_i|} \sum_{t=1}^{|y_i|} 
s_{i,t}(\theta) \widehat{A}_{i,t}
\right] \\
=&\ 
\mathbb{E}_{ x \sim \mathcal{D},\, \{y_i\}_{i=1}^G \sim \pi_{\theta_\text{old}}( \cdot | x) }
\left[ 
\frac{1}{G} \sum_{i=1}^{G} s_i (\theta) \cdot \frac{1}{|y_i|} \sum_{t=1}^{|y_i|} 
\widehat{A}_{i,t} \frac{ \nabla_{\theta} \pi_{\theta} (y_{i,t} | x, y_{i,<t}) }{ \pi_{\theta} (y_{i,t} | x, y_{i,<t}) }
\right] \\
=&\ 
\mathbb{E}_{ x \sim \mathcal{D},\, \{y_i\}_{i=1}^G \sim \pi_{\theta_\text{old}}( \cdot | x) }
\left[ 
\frac{1}{G} \sum_{i=1}^{G}  \left( \frac{ \pi_{\theta} (y_i | x) }{ \pi_{\theta_\text{old}} (y_i | x)} \right)^{\frac{1}{|y_i|}}
\cdot \frac{1}{|y_i|} \sum_{t=1}^{|y_i|}
\widehat{A}_{i,t} \nabla_{\theta} \log \pi_{\theta} (y_{i,t} | x, y_{i,<t})  
\right].
\label{equ:gspo-token_grad}
\end{align}

Observe that the expression $\frac{ \pi_{\theta} (y_{i,t} | x, y_{i,<t}) }{ \mathrm{sg} \left[ \pi_{\theta} (y_{i,t} | x, y_{i,<t}) \right] }$ evaluates to 1, implying that $s_{i,t}(\theta)$ and $s_{i}(\theta)$ yield the same numerical result. By inspecting Equation~\eqref{equ:gspo} versus~\eqref{equ:gspo-token}, as well as Equation~\eqref{equ:gspo_grad} and~\eqref{equ:gspo-token_grad}, one can see that GSPO-token and GSPO coincide numerically in terms of their optimization objective, clipping rule, and theoretical gradient, provided that the advantage values assigned to each token in the sequence $y_i$ are held constant (i.e., $\widehat{A}_{i,t} = \widehat{A}_{i}$). However, GSPO-token offers additional adaptability by allowing token-specific advantages.

\section{Experiments and Discussion}

\subsection{Empirical Results}

We conduct our experiments using a cold-start model that has been fine-tuned from Qwen3-30B-A3B-Base, and we present both the training reward trajectories and evaluation curves for model performance on three benchmarks: AIME'24 (measured as mean Pass@1 over 32 samples), LiveCodeBench (from 202410 to 202502, assessed as average Pass@1 with 8 samples), and CodeForces (using Elo Rating). Throughout the RL training process, each rollout batch is divided into four separate mini-batches for the purposes of gradient updating. For GSPO, we designate the left and right clip intervals in Equation~\eqref{equ:gspo} as 3e-4 and 4e-4, in that order. For a baseline comparison, we use GRPO, fixing the corresponding clipping limits in Equation~\eqref{equ:grpo} to 0.2 (left) and 0.27 (right); these hyperparameters are tuned to facilitate an equitable evaluation. It is important to highlight that GRPO requires the Routing Replay training technique to achieve reliable MoE RL convergence—a point elaborated in \S~\ref{subsec:routing_replay}. By contrast, \textit{GSPO eliminates the necessity for this approach}.

As illustrated in Figure~\ref{fig:results}, GSPO exhibits stable progress during the training process. Notably, \textbf{GSPO consistently enhances model performance by scaling up training computation, periodically refreshing the query set, and increasing generation length}. Additionally, GSPO outperforms GRPO in terms of training efficiency, securing higher training accuracy and improved benchmark results while utilizing equivalent computational resources and queries. Furthermore, the successful implementation of GSPO in the RL training of the recent Qwen3 models robustly demonstrates GSPO’s effectiveness in leveraging RL scaling for large language models.

\subsection{Curious Observation on Clipping Fractions}

An important differentiator between GSPO and GRPO lies in GSPO's approach of clipping at the response level, as opposed to the token level utilized by GRPO. As depicted in Figure~\ref{fig:clipping}, there exists a disparity of approximately two orders of magnitude in the proportion of clipped tokens when comparing GSPO to GRPO (and modifications to the clipping range do not change this fundamental gap). Nonetheless, GSPO attains greater training efficiency than GRPO, even though it discards a substantially larger proportion of tokens for training or gradient estimation purposes. This seemingly paradoxical result—that a method which excludes many more tokens for gradient computation outperforms one that leverages more tokens—underscores the inefficiency and high variance of GRPO's token-wise gradient estimation. On the other hand, GSPO's full-sequence strategy yields a more stable and informative learning signal, enabling improved exploitation of training samples.

\subsection{Benefit of GSPO for MoE Training}
\label{subsec:routing_replay}

\paragraph{Background}
In contrast to reinforcement learning (RL) training with dense models, the inherently sparse activations characteristic of MoE models pose distinctive issues related to training stability. Our investigation revealed that, specifically when employing the GRPO algorithm, the pronounced \textit{expert-activation volatility} in MoE models can obstruct the successful convergence of RL optimization. More precisely, we observed that following one or more rounds of gradient updates, the specific experts chosen to process the same output may shift considerably. As an illustrative case, for the 48-layer Qwen3-30B-A3B-Base configuration, each RL-driven gradient step results in approximately 10\% of the experts engaged under the updated policy $\pi_{\theta}$ differing from those under the previous policy $\pi_{\theta_\text{old}}$ for any given rollout instance. This variability intensifies as the MoE architecture deepens, resulting in token-level importance weights $w_{i,t}(\theta) = \frac{ \pi_{\theta} (y_{i,t} | x, y_{i,<t}) }{ \pi_{\theta_\text{old}} (y_{i,t} | x,y_{i,<t})}$ exhibiting significant fluctuations. Such erratic behavior undermines the reliability of these ratios, as previously examined in \S~\ref{sec:motivation} and~\ref{subsec:gradient}, and disrupts stable RL convergence.

\paragraph{Our Previous Approach} In earlier work, we addressed this problem using the \textbf{Routing Replay} training method. More specifically, this involved storing the experts selected (i.e., the activated experts) under $\pi_{\theta_\text{old}}$ and subsequently ``replaying'' these same expert choices in $\pi_{\theta}$ during the calculation of the importance weights $w_{i,t}(\theta) = \frac{ \pi_{\theta} (y_{i,t} | x, y_{i,<t}) }{ \pi_{\theta_\text{old}} (y_{i,t} | x,y_{i,<t})}$. Consequently, for any given token $y_{i,t}$, the conditional probabilities $\pi_{\theta} (y_{i,t} | x, y_{i,<t})$ and $\pi_{\theta_\text{old}} (y_{i,t} | x,y_{i,<t})$ are computed based on identical expert subnetworks. This approach stabilizes the per-token importance ratios and maintains consistency in which experts are optimized throughout successive gradient updates. As depicted in Figure~\ref{fig:routing_replay}, Routing Replay emerges as a crucial component in facilitating the stable convergence of MoE model training with GRPO.

\paragraph{Benefit of GSPO} While Routing Replay allows GRPO training to achieve successful convergence in MoE models, its mechanism of repeatedly using routing states leads to increased memory and communication demands, and may impose constraints on the usable capacity of the MoE architecture. On the other hand, as depicted in Figure~\ref{fig:results}, GSPO removes reliance on Routing Replay by enabling standard computation of the importance ratios $s_i(\theta)$, ensuring normal convergence behavior and robust optimization throughout training. The central observation here is that GSPO exclusively considers the sequence-level likelihood, $\pi_{\theta} (y_i | x)$, making it insensitive to token-level probabilities such as $\pi_{\theta} (y_{i,t} | x, y_{i,<t})$. Given that the MoE model continuously upholds its core language modeling abilities, the sequence-level likelihood remains relatively stable. Thus, GSPO inherently overcomes the problem of volatility in expert activation within MoE models, eliminating the necessity for intricate mechanisms like Routing Replay. This leads to not only a more streamlined and reliable training procedure, but also enables the model to fully utilize its capacity free from artificial limitations.

\subsection{Benefit of GSPO for RL Infrastructure}

Due to precision inconsistencies that arise between training engines (such as Megatron) and inference engines (including SGLang and vLLM), the standard approach has been to rely on the training engine to re-evaluate the likelihoods of generated responses according to the previous policy $\pi_{\theta_\text{old}}$. In contrast, GSPO employs sequence-level likelihoods during optimization rather than token-level likelihoods, which, by nature, are generally more robust to precision mismatches. Consequently, GSPO enables the direct utilization of likelihoods produced by the inference engine for optimization, removing the necessity for recalculation using the training engine. This is especially advantageous in contexts such as partial rollout, multi-turn RL, or architectures that separate training and inference.

\section{Conclusion}

This work introduces Group Sequence Policy Optimization (GSPO), a novel reinforcement learning algorithm tailored for large language model training. GSPO employs sequence-level importance ratios rooted in importance sampling theory, and incorporates clipping, reward assignment, and optimization directly at the sequence level. Empirical results show that GSPO outperforms GRPO in terms of training robustness, sample efficiency, and overall effectiveness, especially when applied to large-scale reinforcement learning on MoE architectures. These advantages have contributed substantially to the outstanding progress observed in the latest Qwen3 models. GSPO thus provides a robust and extensible algorithmic base for continued scaling of reinforcement learning, which is anticipated to drive significant advancements in artificial intelligence.

\end{document}